\subsection{Khai thác tương quan giao thông và xác định tuyến mục tiêu}
\label{subsec:correlation_target_selection}

Sau khi xây dựng vùng giao thông chức năng như trình bày ở Mục~\ref{subsec:dynamic_zone_construction},
bước tiếp theo nhằm làm rõ các mối quan hệ ảnh hưởng giữa các tuyến đường trong vùng và xác định
tập các tuyến mục tiêu cần dự đoán. Trong bối cảnh giao thông đô thị, ùn tắc không lan truyền
một cách ngẫu nhiên mà thường tuân theo các mối quan hệ có hướng và có độ trễ theo thời gian,
phản ánh động lực di chuyển của phương tiện trong mạng lưới giao thông. Do đó, việc khai thác
các mối tương quan giao thông có xét đến độ trễ đóng vai trò then chốt trong việc xác định
các tuyến có ảnh hưởng quan trọng.

\subsubsection{Tương quan chéo có độ trễ trong vùng giao thông}
\label{subsubsec:zone_correlation}

Đối với mỗi vùng giao thông chức năng, nhóm chỉ xét các tuyến đường thuộc vùng này khi
tính toán tương quan, nhằm đảm bảo tính cục bộ và giảm độ phức tạp tính toán. Với mỗi cặp
tuyến đường $(i,j)$ trong vùng, từ chuỗi thời gian vận tốc đã được chuẩn hóa $S_i(t)$ và
$S_j(t)$, hàm tương quan chéo có độ trễ $X_{i,j}(\tau)$ được tính trên một cửa sổ thời gian
lịch sử, với $\tau \in [-\tau_{\max}, \tau_{\max}]$.

Từ hàm tương quan chéo này, nhóm xác định hai đại lượng đặc trưng cho mối quan hệ giữa
hai tuyến đường: (i) trọng số liên kết dương $W^{\text{pos}}_{i,j}$, phản ánh mức độ nổi trội
của đỉnh tương quan so với nền nhiễu, và (ii) độ trễ tương ứng $\tau^{\text{pos}}_{i,j}$,
xác định hướng và thời điểm ảnh hưởng mạnh nhất. Các đại lượng này được tính theo khuôn khổ
mạng tương quan giao thông được đề xuất bởi Guo et al.~\cite{guo2019identifying}.

Việc sử dụng $W^{\text{pos}}_{i,j}$ thay vì giá trị tương quan cực đại thuần túy cho phép
loại bỏ các mối tương quan yếu hoặc ngẫu nhiên, đồng thời làm nổi bật các liên kết có ý nghĩa
thống kê trong động lực giao thông. Trong khi đó, $\tau^{\text{pos}}_{i,j}$ cung cấp thông
tin về quan hệ lead--lag giữa các tuyến, phản ánh hướng lan truyền tiềm năng của ùn tắc.

\subsubsection{Lọc liên kết tương quan theo tiêu chí thống kê và không gian}
\label{subsubsec:link_filtering}

Không phải mọi cặp tuyến có tương quan dương đều mang ý nghĩa chức năng trong bối cảnh lan
truyền ùn tắc. Do đó, nhóm áp dụng một tập các tiêu chí lọc để chỉ giữ lại các liên kết
tương quan mạnh và hợp lý về mặt không gian--thời gian. Cụ thể, một liên kết $(i,j)$ được giữ
lại nếu thỏa mãn đồng thời các điều kiện sau:
\begin{equation}
W^{\text{pos}}_{i,j} \geq W_{\min}, \quad
D_{\min} \leq D_{i,j} \leq D_{\max}, \quad
|\tau^{\text{pos}}_{i,j}| \leq \tau_{\text{cut}},
\end{equation}
trong đó $D_{i,j}$ là khoảng cách Haversine giữa hai tuyến đường $i$ và $j$. Các ngưỡng
$W_{\min}$, $D_{\min}$, $D_{\max}$ và $\tau_{\text{cut}}$ được lựa chọn nhằm cân bằng giữa
việc giữ lại các liên kết có ý nghĩa và loại bỏ nhiễu.

Bước lọc này đảm bảo rằng các liên kết tương quan được sử dụng trong mô hình không chỉ mạnh
về mặt thống kê, mà còn phù hợp với ràng buộc vật lý và động học của giao thông đô thị. Các
liên kết có khoảng cách quá xa hoặc độ trễ quá lớn thường khó phản ánh mối quan hệ ảnh hưởng
trực tiếp và do đó bị loại bỏ.

\subsubsection{Xác định tập tuyến mục tiêu}
\label{subsubsec:target_mask}

Từ tập các liên kết tương quan thỏa mãn các tiêu chí lọc, nhóm xác định \textit{tập các
tuyến mục tiêu} trong mỗi vùng giao thông. Cụ thể, một tuyến đường $j$ được xem là tuyến
mục tiêu nếu nó xuất hiện trong ít nhất một liên kết tương quan được giữ lại với bất kỳ tuyến
nào khác trong vùng. Trong trường hợp số lượng tuyến mục tiêu quá lớn, nhóm chỉ giữ lại
một tập con gồm các tuyến có trọng số liên kết $W^{\text{pos}}$ lớn nhất.
